% Options for packages loaded elsewhere
\PassOptionsToPackage{unicode}{hyperref}
\PassOptionsToPackage{hyphens}{url}
%
\documentclass[
]{article}
\usepackage{amsmath,amssymb}
\usepackage{lmodern}
\usepackage{iftex}
\ifPDFTeX
  \usepackage[T1]{fontenc}
  \usepackage[utf8]{inputenc}
  \usepackage{textcomp} % provide euro and other symbols
\else % if luatex or xetex
  \usepackage{unicode-math}
  \defaultfontfeatures{Scale=MatchLowercase}
  \defaultfontfeatures[\rmfamily]{Ligatures=TeX,Scale=1}
\fi
% Use upquote if available, for straight quotes in verbatim environments
\IfFileExists{upquote.sty}{\usepackage{upquote}}{}
\IfFileExists{microtype.sty}{% use microtype if available
  \usepackage[]{microtype}
  \UseMicrotypeSet[protrusion]{basicmath} % disable protrusion for tt fonts
}{}
\makeatletter
\@ifundefined{KOMAClassName}{% if non-KOMA class
  \IfFileExists{parskip.sty}{%
    \usepackage{parskip}
  }{% else
    \setlength{\parindent}{0pt}
    \setlength{\parskip}{6pt plus 2pt minus 1pt}}
}{% if KOMA class
  \KOMAoptions{parskip=half}}
\makeatother
\usepackage{xcolor}
\usepackage{longtable,booktabs,array}
\usepackage{calc} % for calculating minipage widths
% Correct order of tables after \paragraph or \subparagraph
\usepackage{etoolbox}
\makeatletter
\patchcmd\longtable{\par}{\if@noskipsec\mbox{}\fi\par}{}{}
\makeatother
% Allow footnotes in longtable head/foot
\IfFileExists{footnotehyper.sty}{\usepackage{footnotehyper}}{\usepackage{footnote}}
\makesavenoteenv{longtable}
\setlength{\emergencystretch}{3em} % prevent overfull lines
\providecommand{\tightlist}{%
  \setlength{\itemsep}{0pt}\setlength{\parskip}{0pt}}
\setcounter{secnumdepth}{-\maxdimen} % remove section numbering
\ifLuaTeX
  \usepackage{selnolig}  % disable illegal ligatures
\fi
\IfFileExists{bookmark.sty}{\usepackage{bookmark}}{\usepackage{hyperref}}
\IfFileExists{xurl.sty}{\usepackage{xurl}}{} % add URL line breaks if available
\urlstyle{same} % disable monospaced font for URLs
\hypersetup{
  hidelinks,
  pdfcreator={LaTeX via pandoc}}

\author{}
\date{}

\begin{document}

\hypertarget{metodologuxeda}{%
\subsection{3. METODOLOGÍA}\label{metodologuxeda}}

\hypertarget{enfoque-metodoluxf3gico-y-criterios-de-decisiuxf3n}{%
\subsubsection{3.1 Enfoque metodológico y criterios de
decisión}\label{enfoque-metodoluxf3gico-y-criterios-de-decisiuxf3n}}

El desarrollo de \emph{The Show Verse} se ha abordado mediante una
\textbf{metodología iterativa basada en el ciclo de vida ágil}, adaptada
a un proyecto individual. La elección de este enfoque frente a las
metodologías en cascada se justifica por la naturaleza exploratoria del
proyecto: la integración con APIs de terceros implica incertidumbres
técnicas que no siempre se pueden anticipar en una fase de análisis, y
la prioridad entre funcionalidades puede variar a medida que se
comprende mejor el problema.

El proceso de desarrollo se articula en iteraciones de aproximadamente
una semana de duración. Al inicio de cada iteración se seleccionan del
backlog las funcionalidades a abordar, priorizadas según tres
criterios:\\
- \textbf{Dependencia técnica:} las funcionalidades que sirven de base a
otras (autenticación, cliente de API) se abordan primeras.\\
- \textbf{Valor aportado al usuario:} se priorizan las funcionalidades
más visibles y relevantes para la experiencia de uso.\\
- \textbf{Riesgo técnico:} las integraciones más inciertas se adelantan
para conocer cuanto antes sus restricciones.

\hypertarget{criterios-de-integraciuxf3n-de-apis}{%
\paragraph{Criterios de integración de
APIs}\label{criterios-de-integraciuxf3n-de-apis}}

La selección de las APIs externas ha seguido los siguientes criterios:

\begin{itemize}
\tightlist
\item
  \textbf{Datos completos:} TMDb fue seleccionada frente a otras
  alternativas (The Open Movie Database, Wikidata) por ofrecer la mayor
  cobertura de contenido, soporte en diferentes idiomas, galería de
  imágenes y datos de streaming. Dispone de la documentación más
  completa y una comunidad de desarrolladores activa.\\
\item
  \textbf{Capacidad de tracking:} Trakt.tv fue seleccionada como
  plataforma de sincronización por su API pública v2 bien documentada,
  su soporte completo de series a nivel de episodio, la disponibilidad
  de sus endpoints principales para desarrolladores, y por ser la
  plataforma de \emph{tracking} más referenciada por la comunidad de
  desarrolladores independientes.\\
\item
  \textbf{Complementariedad:} OMDb fue seleccionada únicamente para los
  datos de rating de IMDb y Rotten Tomatoes, que TMDb no proporciona
  directamente.
\end{itemize}

\hypertarget{criterios-de-decisiuxf3n-para-la-autenticaciuxf3n}{%
\paragraph{Criterios de decisión para la
autenticación}\label{criterios-de-decisiuxf3n-para-la-autenticaciuxf3n}}

El problema de la autenticación OAuth 2.0 admitía varias opciones de
implementación en Next.js. Se descartó el uso de librerías de
autenticación genéricas como NextAuth.js / Auth.js por las siguientes
razones: - Trakt.tv no es un proveedor OAuth soportado de forma nativa
por estas librerías. - La personalización necesaria (refresco de tokens
específico de Trakt, intercambio de códigos con headers personalizados)
habría requerido más configuración que implementar el flujo
directamente. - La implementación propia ofrece mayor control sobre el
manejo de errores y el almacenamiento de tokens.

Se optó por implementar el flujo Authorization Code de forma manual
mediante \textbf{API Routes de Next.js}, almacenando los tokens en
cookies \texttt{httpOnly} con atributos \texttt{Secure} y
\texttt{SameSite=Lax}. Esta decisión garantiza que los tokens no sean
accesibles desde JavaScript del navegador (protección frente a XSS) y
que no se envíen en peticiones \emph{cross-site} no intencionadas
(protección parcial frente a CSRF).

\hypertarget{criterios-de-diseuxf1o-de-interfaz}{%
\paragraph{Criterios de diseño de
interfaz}\label{criterios-de-diseuxf1o-de-interfaz}}

Las decisiones de diseño visual han seguido los siguientes principios: -
\textbf{Coherencia con el dominio:} una plataforma de contenido
audiovisual compite visualmente con las propias plataformas de
\emph{streaming}, desde el punto de vista del usuario, lo que justifica
una estética similar. - \textbf{Modo oscuro por defecto:} adecuado para
el consumo de contenido multimedia y preferido por la mayoría de
usuarios. - \textbf{Mobile-first:} el consumo de contenido audiovisual
en dispositivos móviles supera al escritorio en grupos demográficos
menores de 35 años según datos de Nielsen (2023).

\hypertarget{criterios-de-estrategia-de-cachuxe9}{%
\paragraph{Criterios de estrategia de
caché}\label{criterios-de-estrategia-de-cachuxe9}}

Para cada tipo de dato gestionado en la aplicación se ha definido una
estrategia de caché en función de la frecuencia de actualización del
dato en origen:

\begin{longtable}[]{@{}
  >{\raggedright\arraybackslash}p{(\columnwidth - 4\tabcolsep) * \real{0.3333}}
  >{\raggedright\arraybackslash}p{(\columnwidth - 4\tabcolsep) * \real{0.3333}}
  >{\raggedright\arraybackslash}p{(\columnwidth - 4\tabcolsep) * \real{0.3333}}@{}}
\toprule()
\begin{minipage}[b]{\linewidth}\raggedright
Tipo de dato
\end{minipage} & \begin{minipage}[b]{\linewidth}\raggedright
Almacenamiento
\end{minipage} & \begin{minipage}[b]{\linewidth}\raggedright
Revalidación
\end{minipage} \\
\midrule()
\endhead
Catálogo general (top rated, populares) & ISR servidor (force-cache) &
10 minutos \\
Detalles de película/serie & ISR servidor (force-cache) & 1 hora \\
Imágenes TMDb & CDN Vercel & 1 hora \\
Estado personal del usuario (favoritos, historial) & No-cache & En cada
petición \\
Datos de autenticación (token) & Cookie httpOnly & Por expiración \\
\textbf{Puntuaciones IMDb y Trakt en listas} & \textbf{localStorage
(cliente)} & \textbf{30 días} \\
\textbf{Datos OMDb en página de detalle} (IMDb rating, RT, Metacritic) &
\textbf{sessionStorage (cliente)} & \textbf{24 horas} \\
\textbf{Rating personal del usuario en página de detalle} &
\textbf{Memoria (Map en módulo)} & \textbf{10 min (nulo: 45 s)} \\
\bottomrule()
\end{longtable}

Esta diferenciación permite reducir el número de llamadas a APIs
externas en un estimado del 70--80\% para los datos de catálogo,
mientras se garantiza la actualización de los datos de usuario. El nivel
de caché en cliente dedicado a las puntuaciones de fuentes externas
(IMDb, Trakt, OMDb) reduce adicionalmente la carga sobre estas APIs en
las vistas de listas y detalles, donde se consultan simultáneamente
puntuaciones de múltiples títulos.

\hypertarget{tecnologuxedas-empleadas}{%
\subsubsection{3.2 Tecnologías
empleadas}\label{tecnologuxedas-empleadas}}

\textbf{Next.js 16 (Framework principal)}\\
Seleccionado frente a otras alternativas (Remix, SvelteKit, Nuxt 3) por
ser el framework React con mayor adopción industrial (\textgreater40\%
de cuota entre frameworks React según npm trends 2024), la mejor
integración con Vercel (plataforma de despliegue objetivo) y la
documentación y comunidad más extensa. El App Router de Next.js 13+ es
la única opción que implementa React Server Components de forma estable
y con soporte oficial.

\textbf{React 19}\\
La elección de React como biblioteca de UI es consecuencia directa de la
elección de Next.js. React 19 introduce mejoras en el manejo de acciones
de servidor y la promesa nativa en componentes de cliente, aunque en
este proyecto se utiliza primariamente en su versión estable de
componentes de cliente e hidratación.

\textbf{Tailwind CSS 4}\\
Seleccionado frente a CSS Modules, Styled Components o Emotion por el
paradigma \emph{utility-first} que acelera el desarrollo sin sacrificar
la personalización. La versión 4, basada en PostCSS con motor Rust
nativo, elimina el \emph{PurgeCSS} separado e incorpora mejoras
sustanciales de rendimiento en el tiempo de compilación.

\textbf{Framer Motion 12}\\
La elección de Framer Motion frente a otras librerías de animación
(React Spring, CSS Animations puras) se basa en su modelo declarativo
basado en variantes, que permite expresar animaciones complejas de forma
legible, su soporte nativo para animaciones de \emph{layout} (el
\emph{Flip Animation Technique}), y su integración con el ciclo de vida
de los componentes React mediante \texttt{AnimatePresence}.

\textbf{TypeScript 5}\\
Configurado en modo de verificación estricta para la configuración del
proyecto y los tipos de datos de las respuestas de API. En los
componentes React se utiliza JavaScript con JSX, manteniendo la
flexibilidad de tipado dinámico en la lógica de componentes.

\begin{center}\rule{0.5\linewidth}{0.5pt}\end{center}

\hypertarget{desarrollo}{%
\subsection{4. DESARROLLO}\label{desarrollo}}

\hypertarget{arquitectura-del-sistema}{%
\subsubsection{4.1 Arquitectura del
sistema}\label{arquitectura-del-sistema}}

El sistema se estructura en tres capas diferenciadas, implementadas en
una única base de código Next.js:

\textbf{Capa de presentación (Client Components)}\\
Componentes React que se ejecutan y se hidratan en el navegador.
Gestionan el estado local de la interfaz, las interacciones del usuario
y las peticiones a la capa intermedia. Los componentes de mayor
complejidad son \texttt{MainDashboardClient.jsx} (gestión del dashboard
completo), \texttt{DetailsClient.jsx} (página de detalle de película y
serie) y \texttt{DiscoverClient.jsx} (módulo de búsqueda).

\textbf{Capa de servidor (Server Components + API Routes)}\\
Los Server Components de Next.js realizan la búsqueda inicial de datos
directamente en el servidor, sin coste de JavaScript para el cliente y
las API Routes actúan como proxy hacia las APIs externas, añadiendo los
tokens de autenticación en el servidor y transformando los datos al
formato que consume el frontend.

\textbf{Capa de datos (APIs externas)}\\
TMDb, Trakt.tv y OMDb como fuentes para los datos de catálogo y del
usuario respectivamente.

La separación entre estas capas se hace explícita en la convención de
nombres del proyecto: los archivos terminados en \texttt{Client.jsx} son
siempre Client Components (marcados con \texttt{"use\ client"}),
mientras que los archivos \texttt{page.jsx} y los de la carpeta
\texttt{api/} son siempre Server Components o API Routes.

\hypertarget{implementaciuxf3n-del-sistema-de-autenticaciuxf3n-oauth-2.0}{%
\subsubsection{4.2 Implementación del sistema de autenticación OAuth
2.0}\label{implementaciuxf3n-del-sistema-de-autenticaciuxf3n-oauth-2.0}}

La implementación del flujo OAuth 2.0 con Trakt.tv es el núcleo técnico
más delicado del proyecto. El flujo se articula en cuatro API Routes de
Next.js:

\textbf{Inicio del flujo (\texttt{/api/trakt/oauth}):} construye la URL
de autorización de Trakt con los parámetros \texttt{client\_id},
\texttt{redirect\_uri} y \texttt{response\_type=code}, y redirige al
usuario. Este endpoint se invoca al pulsar el botón ``Conectar con
Trakt'' de la Navbar.

\textbf{Callback (\texttt{/api/trakt/auth/callback}):} recibe el
\texttt{authorization\_code} de Trakt, lo intercambia por
\texttt{access\_token} y \texttt{refresh\_token} mediante una petición
POST servidor-a-servidor a \texttt{https://api.trakt.tv/oauth/token}, y
almacena ambos tokens en cookies \texttt{httpOnly}. A continuación
redirige al usuario a la página de inicio.

\textbf{Estado (\texttt{/api/trakt/auth/status}):} lee las cookies de
sesión y, si existe \texttt{access\_token} válido, hace una petición a
\texttt{/users/me} de Trakt para obtener el perfil del usuario. Si el
token está expirado, intenta refrescarlo usando el
\texttt{refresh\_token} antes de declarar al usuario como no
autenticado.

\textbf{Desconexión (\texttt{/api/trakt/auth/disconnect}):} revoca el
token en Trakt mediante la API \texttt{/oauth/revoke} y elimina las
cookies de sesión.

Este diseño garantiza que ni el \texttt{client\_secret} de Trakt ni los
\texttt{access\_token} del usuario sean accesibles desde el JavaScript
del navegador en ningún momento del flujo.

\hypertarget{integraciuxf3n-con-las-apis-externas}{%
\subsubsection{4.3 Integración con las APIs
externas}\label{integraciuxf3n-con-las-apis-externas}}

\hypertarget{cliente-tmdb}{%
\paragraph{Cliente TMDb}\label{cliente-tmdb}}

El cliente de TMDb (\texttt{tmdb.js}) encapsula toda la lógica de acceso
a la API v3. La función central \texttt{tmdb(path,\ params,\ options)}
implementa:

\begin{itemize}
\tightlist
\item
  Un \textbf{constructor de URL unificado} que añade automáticamente la
  API Key y el idioma por defecto (\texttt{es-ES}).
\item
  Un \textbf{mecanismo de timeout} mediante \texttt{AbortController} (8
  segundos por defecto en cliente, ajustable por opción) para evitar que
  peticiones lentas bloqueen la interfaz.
\item
  \textbf{Comportamiento de caché diferenciado} según el entorno:
  \texttt{force-cache} con revalidación ISR en servidor,
  \texttt{no-store} en cliente (el navegador gestiona su propia caché
  HTTP).
\item
  \textbf{Gestión granular de errores:} los errores 404 retornan
  \texttt{null} silenciosamente (contenido no encontrado es un caso
  esperado); otros errores se loguean y también retornan \texttt{null}
  para que el componente gestione el estado de error.
\end{itemize}

Sobre esta función base se construyen más de 40 funciones nombradas por
dominio (\texttt{fetchTopRatedMovies}, \texttt{getDetails},
\texttt{getCredits}, \texttt{getWatchProviders},
\texttt{getActorDetails}, etc.) que constituyen la API pública del
módulo.

Un aspecto técnico importante es la resolución de imágenes: la función
\texttt{getLogos} obtiene el logo oficial de un título (endpoint
\texttt{/images} de TMDb) aplicando un criterio de selección: prioriza
logos en español, luego en inglés, luego sin idioma; dentro de cada
grupo selecciona el de mayor número de votos. Este logo se usa como
título visual en la página de detalle cuando está disponible, en lugar
del texto plano del título.

\hypertarget{proxy-de-trakt-y-transformaciuxf3n-de-datos}{%
\paragraph{Proxy de Trakt y transformación de
datos}\label{proxy-de-trakt-y-transformaciuxf3n-de-datos}}

Las API Routes de Trakt actúan como un intermediario que resuelve dos
problemas estructurales:

\textbf{Problema de identificadores:} el frontend trabaja exclusivamente
con \texttt{tmdbId} (identificadores de TMDb), mientras que muchos
endpoints de Trakt requieren sus propios identificadores
(\texttt{traktId}, \texttt{slug}). Las API Routes realizan internamente
la traducción mediante el endpoint de búsqueda por ID externo de Trakt
(\texttt{/search/tmdb/\{id\}?type=movie\textbar{}show}).

\textbf{Problema de enriquecimiento de datos:} las respuestas de Trakt
incluyen datos de \emph{tracking} pero no imágenes. Las API Routes que
devuelven listas de contenido (historial, favoritos, pendientes)
realizan peticiones paralelas a TMDb para obtener las URLs de poster y
backdrop, enriqueciendo la respuesta antes de enviarla al cliente.

El cliente del frontend (\texttt{traktClient.js}) expone una API
orientada a operaciones de usuario: \texttt{traktSetWatched},
\texttt{traktSetEpisodeWatched}, \texttt{traktHistoryOp},
\texttt{traktSetRating}, etc. Internamente, todas estas funciones
realizan peticiones fetch a las API Routes internas, que son las que
contienen el token de autenticación y la lógica de llamada a Trakt.

Una particularidad técnica de este módulo es el manejo de fechas: Trakt
acepta formatos distintos según el endpoint (ISO 8601 completo para el
historial de \emph{plays}; \texttt{YYYY-MM-DD} para favoritos y
\emph{watchlist}). El cliente incluye las funciones
\texttt{normalizeWatchedAtForHistoryApi} y
\texttt{normalizeWatchedAtForApi} que normalizan cualquier entrada
(objeto \texttt{Date}, string en varios formatos) al formato correcto
para cada caso.

\hypertarget{cachuxe9-de-puntuaciones-en-listas-de-favoritos-y-pendientes}{%
\paragraph{Caché de puntuaciones en listas de favoritos y
pendientes}\label{cachuxe9-de-puntuaciones-en-listas-de-favoritos-y-pendientes}}

Las páginas de Favoritos (\texttt{FavoritesClient.jsx}) y
Pendientes(\texttt{WatchlistClient.jsx}) muestran, por cada ítem de la
lista, las puntuaciones IMDb y Trakt de la comunidad junto con la
valoración personal del usuario en IMDb y Trakt. Estas puntuaciones se
obtienen de las APIs externas (OMDb, Trakt \emph{scoreboard} endpoint) y
deben cargarse de forma asíncrona y paginada para no exceder los límites
de peticiones.

Para evitar peticiones redundantes entre sesiones y navegaciones, se
implementa un \textbf{sistema de caché de puntuaciones en tres niveles}:

\textbf{Nivel 1 --- \texttt{localStorage} con TTL de 30 días
(puntuaciones IMDb y Trakt).} Las funciones
\texttt{readScoreCache(source)}, \texttt{writeScoreCache(source,\ map)}
y \texttt{updateScoreCache(source,\ id,\ score)} gestionan un
diccionario persistente en \texttt{localStorage} bajo las claves
\texttt{showverse:scores:imdb} y \texttt{showverse:scores:trakt}. Cada
entrada almacena el valor numérico y el \emph{timestamp} de escritura
(\texttt{\{\ score,\ t\ \}}). En la lectura, las entradas con antigüedad
superior a 30 días se descartan automáticamente. Este TTL largo está
justificado porque las puntuaciones de comunidad de IMDb y Trakt cambian
con muy baja frecuencia.

\textbf{Nivel 2 --- \texttt{sessionStorage} con TTL de 24 horas (datos
OMDb).} El módulo \texttt{omdbCache.js} gestiona una caché en
\texttt{sessionStorage} bajo la clave
\texttt{showverse:omdb:\{imdbId\}}, que almacena el rating de IMDb, los
votos, los premios y las puntuaciones de Rotten Tomatoes y Metacritic
obtenidos de la API de OMDb. Este mismo módulo se reutiliza en la página
de detalle y en la lista de favoritos para evitar llamadas duplicadas a
OMDb durante una misma sesión de navegación.

\textbf{Nivel 3 --- Caché en memoria (\texttt{Map}) con TTL corto
(rating personal del usuario).} Los mapas \texttt{userRatingCache} y
\texttt{traktScoreCache} (instancias de \texttt{Map} a nivel de módulo)
almacenan en memoria las valoraciones personales del usuario autenticado
y las puntuaciones de comunidad de Trakt, con TTL de 10 minutos para
entradas con valor y 45 segundos para entradas nulas (para reintentar si
la API no respondió). Estos TTL cortos garantizan que los datos de
usuario se actualizan con frecuencia razonable, mientras se evitan
peticiones en cada \emph{re-render}.

\textbf{Carga diferida y patrón \emph{on-hover}.} Para minimizar el
número de peticiones en el montaje inicial del componente, las
puntuaciones IMDb y Trakt se cargan de forma diferida al primer
\emph{hover} sobre una tarjeta de la lista (\texttt{handleHover}). Si el
valor ya está en caché de nivel 1, se recupera instantáneamente sin
llamada de red. La función \texttt{runPool(items,\ limit,\ worker)}
controla la concurrencia máxima de peticiones simultáneas para respetar
los \emph{rate limits} de las APIs externas.

\hypertarget{barra-de-puntuaciones-en-puxe1gina-de-detalle-scoreboardbar}{%
\paragraph{\texorpdfstring{Barra de puntuaciones en página de detalle
(\texttt{ScoreboardBar})}{Barra de puntuaciones en página de detalle (ScoreboardBar)}}\label{barra-de-puntuaciones-en-puxe1gina-de-detalle-scoreboardbar}}

El componente \texttt{ScoreboardBar.jsx} centraliza la visualización de
puntuaciones en la página de detalle de película o serie: puntuación
TMDb (procedente de los datos ya cargados por el Server Component, sin
petición adicional), puntuación de comunidad de Trakt (obtenida del
endpoint \texttt{/api/trakt/scoreboard}) y rating de IMDb (procedente de
OMDb a través de \texttt{omdbCache}). El componente también gestiona la
valoración personal del usuario en Trakt (estrellas 1--10), con
actualización optimista del estado local.

\hypertarget{implementaciuxf3n-de-los-muxf3dulos-funcionales-principales}{%
\subsubsection{4.4 Implementación de los módulos funcionales
principales}\label{implementaciuxf3n-de-los-muxf3dulos-funcionales-principales}}

\hypertarget{dashboard-principal}{%
\paragraph{Dashboard principal}\label{dashboard-principal}}

La página de inicio (\texttt{/}) es un Server Component que obtiene en
paralelo los datos de todas las secciones mediante \texttt{Promise.all},
reduciendo el tiempo total de carga a la latencia de la llamada más
lenta en lugar de la suma de todas ellas. Para el hero, se aplica una
selección inteligente del backdrop: se descartan imágenes sin texto en
inglés (inadecuadas para el hero), se priorizan las de mayor resolución
(1280px) y, dentro de estas, las de mayor número de votos.

El componente cliente \texttt{MainDashboardClient.jsx} recibe todos los
datos del servidor ya preparados y se encarga exclusivamente de la
composición visual y las interacciones: el carrusel automático del hero
(con pausa al hover), el sticky scroll de la Navbar y la aparición
progresiva de las secciones mediante \texttt{staggerChildren} de Framer
Motion.

\hypertarget{seguimiento-por-episodio}{%
\paragraph{Seguimiento por episodio}\label{seguimiento-por-episodio}}

El módulo de episodios es el funcionalmente más complejo. El estado de
visionado por episodio se representa mediante el objeto
\texttt{watchedBySeason}.

Este objeto se obtiene del endpoint \texttt{/api/trakt/show/watched},
que hace una llamada a \texttt{/sync/watched/shows} de Trakt y
transforma la respuesta para indexarla por temporada y número de
episodio. Cada toggle de episodio actualiza este objeto optimistamente
en el estado local del componente, mientras la petición al servidor se
procesa en segundo plano.

El botón ``Marcar temporada completa'' construye la lista de episodios
no vistos de la temporada y realiza una única llamada a Trakt,
reduciendo el número de peticiones de N a 1.

\hypertarget{actualizaciuxf3n-de-la-interfaz}{%
\paragraph{Actualización de la
interfaz}\label{actualizaciuxf3n-de-la-interfaz}}

Para las operaciones de favoritos, pendientes e historial de visionado,
se implementa el siguiente patrón de actualización de la interfaz:

\begin{enumerate}
\def\labelenumi{\arabic{enumi}.}
\tightlist
\item
  Cuando el usuario realiza una acción, el estado local se actualiza
  inmediatamente (el botón cambia de estado, el contador de la Navbar se
  actualiza).
\item
  En paralelo, se envía la petición al servidor.
\item
  Si la petición falla, el estado local se revierte al valor anterior y
  se muestra una notificación de error.
\end{enumerate}

Este patrón elimina el retraso perceptible entre la acción del usuario y
la respuesta de la interfaz, mejorando significativamente la experiencia
de usuario.

\hypertarget{diseuxf1o-de-la-interfaz-y-sistema-de-vistas}{%
\subsubsection{4.5 Diseño de la interfaz y sistema de
vistas}\label{diseuxf1o-de-la-interfaz-y-sistema-de-vistas}}

La interfaz implementa el sistema de diseño descrito en la metodología.
El fondo negro (\texttt{\#000000}) como color base crea el lienzo sobre
el que los fondos (\emph{backdrops}) de las películas y series actúan
como elemento visual protagonista. Las superficies de las tarjetas y los
paneles laterales utilizan valores de opacidad muy bajos
(\texttt{rgba(255,255,255,0.05)}) para el efecto visual, con bordes de
\texttt{1px\ solid\ rgba(255,255,255,0.1)}.

El sistema de tres vistas (Grid, List, Compact) es un componente
transversal que se aplica en varios módulos de la aplicación (Favoritos,
Pendientes, Historial, En Progreso). La selección de vista se persiste
en \texttt{localStorage} por módulo, de forma que el usuario puede tener
configuraciones distintas en cada sección. La transición entre vistas
utiliza \texttt{AnimatePresence} de Framer Motion con un efecto de
\emph{cross-fade} suave, evitando saltos bruscos en el \emph{layout}.

\hypertarget{rendimiento-en-producciuxf3n}{%
\subsubsection{4.6 Rendimiento en
producción}\label{rendimiento-en-producciuxf3n}}

Las métricas obtenidas en producción (Vercel) mediante Lighthouse
validan las decisiones de arquitectura tomadas:

\begin{longtable}[]{@{}lll@{}}
\toprule()
Métrica & Valor & Objetivo \\
\midrule()
\endhead
Performance (Lighthouse) & 92/100 & \textgreater{} 90 \\
SEO (Lighthouse) & 100/100 & 100 \\
Best Practices & 95/100 & \textgreater{} 90 \\
LCP & 1.8s & \textless{} 2.5s \\
CLS & 0.05 & \textless{} 0.1 \\
FID & 45ms & \textless{} 100ms \\
\bottomrule()
\end{longtable}

El LCP de 1.8s se consigue principalmente gracias al SSR: el HTML con el
hero ya renderizado llega al navegador antes de que se descargue todo el
JavaScript, permitiendo que el motor de renderizado del navegador
comience a mostrar contenido de inmediato. El CLS de 0.05 se logra
mediante el uso sistemático de \texttt{width} y \texttt{height}
explícitos en el componente \texttt{\textless{}Image\textgreater{}} de
Next.js, que reserva el espacio de la imagen antes de que esta se
descargue.

\begin{center}\rule{0.5\linewidth}{0.5pt}\end{center}

\hypertarget{conclusiones}{%
\subsection{5. CONCLUSIONES}\label{conclusiones}}

El trabajo ha demostrado que los objetivos planteados son alcanzables
con las tecnologías actuales de desarrollo web y sin necesidad de
infraestructura de servidor propia. Los resultados obtenidos permiten
extraer las siguientes conclusiones:

En primer lugar, la arquitectura Next.js App Router con React Server
Components ha demostrado ser particularmente adecuada para aplicaciones
de este tipo, que combinan una gran cantidad de datos externos con la
necesidad de interactividad en cliente. La separación entre la obtención
de datos (servidor) y la presentación interactiva (cliente) simplifica
el flujo de datos y reduce las fuentes de error. Los Server Components
permiten realizar fetching paralelo de múltiples APIs en el servidor sin
coste para el cliente, lo que resulta en tiempos de carga
significativamente inferiores a los de una SPA equivalente.

En segundo lugar, la integración de OAuth 2.0 mediante API Routes de
Next.js ha probado ser una estrategia robusta y segura sin necesidad de
librerías de autenticación genéricas. El patrón de almacenar tokens en
cookies \texttt{httpOnly} y delegar toda la lógica de autenticación al
servidor cumple con las recomendaciones de seguridad del RFC 6749 y las
guías de OWASP para aplicaciones OAuth.

En tercer lugar, la estrategia de caché diferenciada por tipo de dato ha
demostrado ser eficaz en todas sus capas: ISR en servidor para el
catálogo (servido desde caché en prácticamente el 100\% de las
peticiones), \texttt{no-store} para los datos de usuario protegidos por
autenticación, y una caché en cliente de tres niveles para las
puntuaciones externas de IMDb, Trakt y OMDb en las vistas de listas.
Esta última capa ---\texttt{localStorage} con TTL de 30 días para
puntuaciones de comunidad, \texttt{sessionStorage} con TTL de 24 horas
para datos de OMDb y caché en memoria para el rating personal del
usuario--- reduce drásticamente el número de llamadas a APIs externas en
las páginas de Favoritos y \emph{Watchlist}, donde se consultan
simultáneamente puntuaciones de decenas de títulos.

Finalmente, desde el punto de vista del diseño, el uso de Framer Motion
para las animaciones y el sistema de superficies semi-transparentes con
desenfoque de fondo ha producido una interfaz visualmente diferenciada
respecto a las plataformas de \emph{tracking} existentes, más próxima en
estética a las propias plataformas de \emph{streaming}. Las métricas de
rendimiento obtenidas son completamente compatibles con un diseño
visualmente rico, siempre que las animaciones se implementen sobre la
GPU (propiedades \texttt{transform} y \texttt{opacity}) y no fuercen la
recarga del DOM.

\hypertarget{discusiuxf3n}{%
\subsubsection{5.1 Discusión}\label{discusiuxf3n}}

El trabajo presenta algunas limitaciones que deben reconocerse:

\textbf{Ausencia de tests automatizados.} El proyecto no incluyó una
suite de tests unitarios ni de integración durante el desarrollo. La
validación se realizó exclusivamente mediante pruebas manuales y esta
decisión, justificada por las restricciones temporales del TFG, supone
el principal déficit de calidad del proyecto desde el punto de vista del
proceso de ingeniería del software. Los componentes con lógica de
negocio compleja (cálculo de progreso de episodios, normalización de
fechas para Trakt, estrategia de selección de backdrops) serían
candidatos prioritarios para la cobertura de tests unitarios.

\textbf{Dependencia de APIs de terceros.} La plataforma depende de la
continuidad y estabilidad de tres APIs externas (TMDb, Trakt.tv, OMDb),
ninguna de las cuales está bajo el control del proyecto. Cambios en sus
modelos de precios, modificaciones de sus límites de peticiones o
interrupciones de servicio afectarían directamente a la aplicación. Esta
es una limitación estructural de la arquitectura elegida, aceptada
conscientemente por las ventajas que ofrece frente a la alternativa de
construir una base de datos propia.

\textbf{Accesibilidad parcial.} La puntuación de accesibilidad en
Lighthouse (88/100) indica que, aunque se han aplicado prácticas básicas
(textos alternativos en imágenes, contraste de colores suficiente), no
se ha alcanzado el nivel WCAG 2.1 AA completo. En particular, la
navegación exclusiva por teclado en los modales de episodios y el manejo
del foco en las transiciones de página requieren trabajo adicional.

\textbf{Ausencia de gestión de usuarios propia.} La plataforma delega
toda la gestión de identidad en IMDb y Trakt.tv, lo que significa que
los usuarios deben tener una cuenta en Trakt para acceder a las
funciones personalizadas. Aunque esto reduce la complejidad de
implementación, limita la autonomía de la plataforma.

\hypertarget{luxedneas-de-investigaciuxf3n-futuras}{%
\subsubsection{5.2 Líneas de investigación
futuras}\label{luxedneas-de-investigaciuxf3n-futuras}}

A partir del trabajo realizado, se identifican las siguientes líneas de
desarrollo futuro:

\textbf{Implementación de Progressive Web App (PWA).} La conversión de
la aplicación en una PWA, mediante la adición de un \emph{Service
Worker} y un \emph{manifest} web, permitiría la instalación en
dispositivos móviles y escritorio, y habilitaría un modo de
funcionamiento \emph{offline} básico (visualización del historial en
caché). Esta línea es de especial interés dado que el consumo de
contenido audiovisual está fuertemente correlacionado con el dispositivo
móvil.

\textbf{Sistema de recomendaciones personalizado.} El historial de
visionado acumulado por el usuario constituye un conjunto de datos con
potencial para la implementación de algoritmos de recomendación
colaborativa o basada en contenido. La integración de un modelo de
\emph{machine learning} ligero (posiblemente mediante TensorFlow.js en
el cliente o una función \emph{serverless}) para enriquecer las
secciones del dashboard con recomendaciones verdaderamente
personalizadas es una extensión natural del trabajo actual.

\textbf{Expansión de las funcionalidades sociales.} La comparación del
historial entre usuarios, las recomendaciones de amigos y la
visualización del perfil público de otros usuarios de Trakt son
funcionalidades que la API de Trakt soporta pero que no han sido
implementadas en la versión actual.

\textbf{Backend propio y base de datos local.} A largo plazo, la
eliminación de la dependencia total de Trakt.tv como fuente de verdad
para los datos del usuario requeriría la implementación de un backend
propio con base de datos (PostgreSQL o MySQL) que almacenara el
historial y las preferencias localmente. Esto permitiría funcionalidades
no disponibles en Trakt (estadísticas avanzadas, exportación de datos en
formatos personalizados) y eliminaría la dependencia de servicio
externo.

\textbf{Integración con plataformas de streaming.} Las APIs de Netflix,
Disney+ y Amazon Prime no son públicas, pero servicios como JustWatch
disponen de datos actualizados sobre disponibilidad de contenido en cada
plataforma y país. La integración de esta información en la página de
detalle permitiría completar el ciclo de usuario: ``¿Dónde puedo ver
este contenido ahora mismo?''.

\begin{center}\rule{0.5\linewidth}{0.5pt}\end{center}

\hypertarget{referencias}{%
\subsection{6. REFERENCIAS}\label{referencias}}

\hypertarget{bibliografuxeda}{%
\subsubsection{6.1 Bibliografía}\label{bibliografuxeda}}

Babich, N. (2021). \emph{Glassmorphism in User Interfaces}. UX Planet.
https://uxplanet.org/glassmorphism-in-user-interfaces-1f39bb1308c9

Berners-Lee, T., Fielding, R., y Masinter, L. (1999). \emph{RFC 2396:
Uniform Resource Identifiers (URI): Generic Syntax}. IETF.
https://www.rfc-editor.org/rfc/rfc2396

Bloch, J. (2008). \emph{Effective Java} (2.ª ed.). Addison-Wesley.

CNMC. (2024). \emph{Informe sobre el consumo de vídeo en España}.
Comisión Nacional de Mercados y la Competencia. https://www.cnmc.es/

Fielding, R. T. (2000). \emph{Architectural Styles and the Design of
Network-based Software Architectures} {[}Tesis doctoral, University of
California, Irvine{]}.
https://www.ics.uci.edu/\textasciitilde fielding/pubs/dissertation/top.htm

Framer. (2024). \emph{Framer Motion Documentation}.
https://www.framer.com/motion/

Google Developers. (2024). \emph{Web Vitals --- Essential metrics for a
healthy site}. https://web.dev/vitals/

Google Developers. (2024). \emph{Core Web Vitals report}.
https://support.google.com/webmasters/answer/9205520

Hardt, D. (Ed.). (2012). \emph{RFC 6749: The OAuth 2.0 Authorization
Framework}. IETF. https://www.rfc-editor.org/rfc/rfc6749

Letterboxd. (2024). \emph{About Letterboxd}.
https://letterboxd.com/about/

Lotz, A. D. (2022). \emph{Netflix and Streaming Video: The Business of
Subscriber-Funded Video on Demand}. Polity Press.

Meta Open Source. (2024). \emph{React Documentation --- React Server
Components}. https://react.dev/reference/rsc/server-components

Nielsen. (2023). \emph{The Nielsen Streaming Unwrapped: 2023 Mid-Year
Report}. Nielsen.
https://www.nielsen.com/insights/2023/streaming-unwrapped/

OASIs. (2019). \emph{OpenAPI Specification 3.0}.
https://spec.openapis.org/oas/v3.0.0

OWASP. (2024). \emph{OWASP Top 10: Broken Authentication}. Open Web
Application Security Project. https://owasp.org/www-project-top-ten/

Schwartz, B. (2004). \emph{The Paradox of Choice: Why More Is Less}.
Ecco/HarperCollins.

Statista. (2024). \emph{Number of Netflix subscribers worldwide from 1st
quarter 2013 to 4th quarter 2023}.
https://www.statista.com/statistics/250934/quarterly-number-of-netflix-streaming-subscribers-worldwide/

The Movie Database. (2024). \emph{TMDb API Documentation v3}.
https://developer.themoviedb.org/docs/getting-started

Trakt. (2024). \emph{Trakt API Documentation v2}.
https://trakt.docs.apiary.io/

Vercel. (2024). \emph{Next.js Documentation --- App Router}.
https://nextjs.org/docs/app

Vercel. (2024). \emph{Vercel Documentation --- Edge Functions and
Middleware}. https://vercel.com/docs/functions

W3C. (2023). \emph{Web Content Accessibility Guidelines (WCAG) 2.1}.
World Wide Web Consortium. https://www.w3.org/TR/WCAG21/

Wieruch, R. (2023). \emph{The Road to Next.js}. Robin Wieruch.
https://www.roadtonextjs.com/

\end{document}
